\documentclass[11pt,a4paper]{article}
\usepackage[T1]{fontenc}
\usepackage{lmodern,microtype,amsmath,amssymb}
\usepackage[margin=25mm]{geometry}
\usepackage[colorlinks=true,urlcolor=blue!60!black,linkcolor=blue!60!black]{hyperref}
\usepackage{xcolor}
\hypersetup{pdftitle={From fidelity monotonicity to von Neumann monotonicity at every order}}
\newcommand{\Tr}{\operatorname{Tr}}
\newcommand{\A}{\mathcal A}
\newcommand{\id}{I}
\setlength{\parindent}{0pt}
\setlength{\parskip}{6pt plus 1pt minus 1pt}
\raggedbottom
\begin{document}
\begin{center}
{\Large\bfseries From fidelity monotonicity to von Neumann\\[3pt] monotonicity at every order}\par
\smallskip
{\normalsize Total anticoherence under $\mathrm{SU}(2)$-covariant channels}
\end{center}

\textbf{Theorem.} Fix $1\le t\le N$ and a completely positive, trace-preserving
$\mathrm{SU}(2)$-covariant channel $\Phi$ on the spin-$N/2$ space.
If $\A_t^{T,F}(\Phi(\eta))\ge\A_t^{T,F}(\eta)$ for every density matrix $\eta$ on that space, then
\[
           \boxed{\quad \A_t^{T,\mathrm{vN}}(\Phi(\rho))
                     \ge\A_t^{T,\mathrm{vN}}(\rho)
                     \quad\text{for every density matrix }\rho.\quad}
\]
Both quantifiers include pure and rank-deficient states.

\textbf{Proof.} Put $n=N+1$, $d=t+1$, and write $R=R_{N,t}$ for the symmetric $t$-qubit marginal. Let $\sigma=R\rho$ and $\tau=R\Phi(\rho)$. In the manuscript's squared-fidelity convention,
\[
 \A_t^{T,F}(\rho)=\frac{(\Tr\sqrt\sigma)^2-1}{t},\qquad
 \A_t^{T,\mathrm{vN}}(\rho)=\frac{S(\sigma)}{\log d},\qquad
 S(X)=-\sum_i\lambda_i(X)\log\lambda_i(X),
\]
with natural logarithms and $0\log 0=0$.

Let $\omega=\id_n/n$. Covariance makes $\Phi(\omega)$ invariant under spin rotations. Irreducibility and trace one give $\Phi(\omega)=\omega$. The symmetric partial trace preserves rotation covariance and trace, and the spin-$t/2$ representation is irreducible. Hence $R\omega=\id_d/d$. For every $s\ge0$ the density matrix
\[
 \rho_s=\frac{\rho+ds\,\omega}{1+ds}
 \quad\text{has marginals}\quad
 R\rho_s=\frac{\sigma+s\id_d}{1+ds},\qquad
 R\Phi(\rho_s)=\frac{\tau+s\id_d}{1+ds}.
\]
Apply the fidelity hypothesis to $\rho_s$. Both square-root traces are nonnegative. Comparing their squares therefore compares the traces themselves, and the common normalization cancels. This gives
\[
      \Delta(s):=\Tr\sqrt{\tau+s\id_d}-\Tr\sqrt{\sigma+s\id_d}\ge0
      \qquad(s\ge0).
\]
The spectral identity proved on the next page is
\begin{equation}
 S(\tau)-S(\sigma)=4\lim_{b\to\infty}\int_0^b\Delta(u^2)\,du
                 =2\int_0^\infty\Delta(s)\,\frac{ds}{\sqrt{s}} . \label{eq:identity}
\end{equation}
Every finite integral is nonnegative. Taking the limit and dividing by $\log d>0$ proves the theorem. In particular, comparing fidelity along the segment from $\rho$ to $\omega$ suffices for the entropy comparison at $\rho$.\hfill$\square$

\newpage
\section*{The shifted square-root identity}
Let $x_1,\ldots,x_d$ and $y_1,\ldots,y_d$ be the eigenvalues of $\sigma$ and $\tau$.
They are nonnegative and have equal sum. For $x\ge0$ and $u\ge0$, put
\[
 P_x(u)=\tfrac12\bigl[u\sqrt{u^2+x}
                   + x\log(u+\sqrt{u^2+x})\bigr].
\]
For $x=0$ take $P_0(u)=u^2/2$, including at the origin. For every $x\ge0$ this function is continuous on $[0,\infty)$, has derivative $\sqrt{u^2+x}$ for $u>0$, and satisfies $P_x(0)=x\log(x)/4$. Thus, for $b\ge0$,
\[
 \int_0^b\sqrt{u^2+x}\,du=P_x(b)-\tfrac14 x\log x.
\]

To evaluate the limit, take $b>0$ and write $v=\sqrt{1+x/b^2}$. Rationalizing the first term and factoring $2b$ inside the logarithm gives the exact identity
\[
 P_x(b)-\tfrac12 b^2-\tfrac x2\log(2b)
       =\frac{x}{2(1+v)}+\frac x2\log\!\left(\frac{1+v}{2}\right)
       \ \longrightarrow\ \tfrac x4 .
\]
Apply the finite integral formula to each eigenvalue and subtract the sum on $x$ from the sum on $y$. Equal dimension cancels the $b^2/2$ terms. Equal trace cancels the $\log(2b)$ terms and the limiting linear terms. What remains is
\[
 4\lim_{b\to\infty}\int_0^b
     \left(\sum_i\sqrt{u^2+y_i}-\sum_i\sqrt{u^2+x_i}\right)du
 =-\sum_i y_i\log y_i+\sum_i x_i\log x_i .
\]
The integrand is continuous at $u=0$, even for zero eigenvalues. Substituting $s=u^2$ on finite intervals and then taking the limit proves the second equality in~\eqref{eq:identity}.

\section*{What the converse would imply at arbitrary order}
Fix $N,t$ as above. Write $F(\Phi)$ and $E(\Phi)$ for the respective
fidelity and entropy inequalities on \emph{all} density matrices at this fixed
order. For the full class of same-spin covariant channels,
\begin{equation}\label{eq:converse}
 \boxed{\quad [\forall\Phi,\ E(\Phi)\Rightarrow F(\Phi)]
       \quad\Longleftrightarrow\quad [\forall\Phi,\ F(\Phi)].\quad}
\end{equation}
In particular, if any channel fails $F$, another channel at the same $N,t$
satisfies $E$ and fails $F$. The existence of a fidelity violation at higher
orders is left open.

\textbf{Proof.} Let $J=I_d/d$, $V(X)=\Tr(X-J)^2$, and suppose $F(\Psi)$ fails.
Interpolate between the completely depolarizing channel $\mathcal D$ and
$\Psi$ by $\Phi_a=(1-a)\mathcal D+a\Psi$, $0\le a\le1$. Each $\Phi_a$ is
covariant and CPTP. For a fixed input, put $\omega=R\Psi(\rho)$ and
$\tau_a=J+a(\omega-J)$. On the irreducible tensor operators of rank $L$,
$\Psi$ acts by a real scalar $f_L$ with $|f_L|\le1$. Indeed it is positive
and unital and hence contracts the operator norm of Hermitian operators.
The symmetric marginal intertwines rotations, so its mutually orthogonal
traceless tensor ranks inherit the same multipliers. Thus
\[
 V(\omega)\le V(\sigma),\qquad V(\tau_a)=a^2 V(\omega).
\]
For each input $T(\tau_a)=\Tr\sqrt{\tau_a}$ is continuous in $a$.
The parameters satisfying $F(\Phi_a)$ form a closed subset of $[0,1]$
containing $0$ (the maximum of $T$ is at $J$) and excluding $1$.
Let $a_*$ be its largest element and put $h=1-a_*^2>0$.

For $s>0$, Taylor's theorem about $1/d$ on $[0,1]$, applied to
$x\mapsto\sqrt{s+x}$, bounds its negative quadratic remainder between
$(x-1/d)^2/[8(s+1)^{3/2}]$ and $(x-1/d)^2/(8s^{3/2})$.
Equal trace cancels the linear terms. At $a_*$ the resulting shifted gap satisfies
\[
 D(s)=T(\tau_{a_*}+sI)-T(\sigma+sI)
 \ \ge\ \frac{V(\sigma)}8
       \bigl[(s+1)^{-3/2}-a_*^2 s^{-3/2}\bigr].
\]
The all-input property $F(\Phi_{a_*})$ also gives $D(s)\ge0$ at every shift.
For $s\ge3/h$, convexity gives $(1+1/s)^{-3/2}\ge1-3/(2s)\ge1-h/2$.
Inserting just this tail into~\eqref{eq:identity} gives a uniform entropy margin
\[
 S(\tau_{a_*})-S(\sigma)
 \ge 2\int_{3/h}^{\infty}\frac{hV(\sigma)}{16s^2}\,ds
 =\frac{h^2}{24}V(\sigma).
\]

Let $y_i$ be the eigenvalues of $\omega-J$ and $c=1/d$.
Since $\sum_i y_i=0$ and $c+a y_i\ge(1-a)c$, the mean value theorem for
$\log$ gives, for $0\le a<1$,
\[
 -\frac{d}{da}S(\tau_a)
   =\sum_i y_i\,[\log(c+a y_i)-\log c]
   \le \frac{da}{1-a}\sum_i y_i^2
   \le \frac{da}{1-a}V(\sigma).
\]
Set $b=a_*+(1-a_*)h^2/(96d)$. On $[a_*,b]$ one has
$1-a\ge(1-a_*)/2$. The loss of entropy up to $b$ is therefore at most
$h^2 V(\sigma)/48$. The same $b$ works for every input, including those
with $V(\sigma)=0$. Consequently $E(\Phi_b)$ holds, whereas
$b>a_*$ gives failure of $F(\Phi_b)$. The reverse implication
in~\eqref{eq:converse} is immediate.\hfill$\square$

The placement of the channel quantifier matters.
Equation~\eqref{eq:converse} addresses $\forall\Phi\,[E(\Phi)\Rightarrow
F(\Phi)]$. Its proof gives no unconditional assertion of universal fidelity
monotonicity from the sole premise $\forall\Phi\,E(\Phi)$ at an unresolved order.

\section*{Lean certificate}
The public theorem \texttt{LowOrder.physical\_vonNeumann\_of\_fidelity} uses
the normalized Dicke marginal, the positive matrix square root and the manuscript's two totals. It holds for $N=k+t$ with $t\ge1$ and $k\ge0$, including the case $N=t$. Its single monotonicity hypothesis is fidelity on all density matrices for the fixed channel. The analytical part checks the finite-cutoff form of~\eqref{eq:identity} for arbitrary nonnegative spectra of equal mass, including zero eigenvalues.
The public theorem for~\eqref{eq:converse} is
\path{LowOrder.Converse.physical_converse_iff_universal_fidelity} in
\path{PhysicalConverse.lean}. Its channel type packages complete positivity,
trace preservation and covariance on density matrices. Its only order
hypothesis is $1\le t$ for $N=k+t$. The Lean proof bounds the entropy
integral on a finite interval of positive shifts and uses a finite
difference estimate for the loss along the depolarizing path. Both give
the strict margin required in the compact-interval argument.
\end{document}
